% ID: SEL-D0073-OA-E03
% ESTADO: cotejada. Segunda pasada visual y borrador compilado revisados.
% Fuente: Selectividad/sel_2008_matematicas_aplicadas/Mat CCSS 1 - Examen y criterios de correción.pdf
% SHA-256: 2461e302e686c60322df8edd0a004b11753e95c4824dc5835b16b76c05174e81
% Página: 1; opción A; ejercicio original 3.

\noindent\underline{Parte I}

Laura tiene en su monedero 6 monedas francesas, 2 italianas y 4 españolas. Vicente tiene 9 francesas y 3 italianas. Cada uno saca, al azar, una moneda de su monedero y observa la nacionalidad.
\begin{enumerate}[label=\alph*)]
\item \textbf{(0.5 puntos)} Obtenga el espacio muestral asociado al experimento.
\item \textbf{(1 punto)} \textquestiondown{}Cuál es la probabilidad de que las monedas extraídas no sean de la misma nacionalidad?
\item \textbf{(0.5 puntos)} \textquestiondown{}Cuál es la probabilidad de que ninguna de las monedas extraídas sea francesa?
\end{enumerate}
\par\noindent\begin{minipage}{\linewidth}
\noindent\underline{Parte II}

Se desea estimar la proporción de individuos zurdos en una determinada ciudad. Para ello se toma una muestra aleatoria de 300 individuos resultando que 45 de ellos son zurdos.
\begin{enumerate}[label=\alph*)]
\item \textbf{(1.5 puntos)} Calcule, usando un nivel de confianza del $97\%$, el correspondiente intervalo de confianza para la proporción de individuos zurdos de la población.
\item \textbf{(0.5 puntos)} \textquestiondown{}Sería mayor o menor el error de estimación si se usara un nivel de confianza del $95\%$? Razone la respuesta.
\end{enumerate}
\end{minipage}\par
